% =====================================================================
% Sections V and VI.  \input{} from main.tex.
% =====================================================================

\section{Discussion and Limitations}
\label{sec:limits}

The result we are most confident in is negative. Context clustering, inserted
into a fixed MetaFormer skeleton and compared against a convolutional branch
at the same parameter count, loses by $1.03$ pp at one to three million
parameters and $32\times32$ input. The supporting machinery built around it
--- fold partitioning, channel shuffle, a tuned occupancy --- moves accuracy
by less than seed variance, and the one setting that matters a great deal,
the stem stride, matters because of resolution rather than because of
clustering. None of this contradicts \cite{coc}, which reports a different
operating point; it bounds where the mechanism is useful.

The positive result is narrower than the headline number suggests, in three
specific ways.

\textbf{It depends on distillation.} Under cross-entropy the proposed model
leads ShuffleNetV2-x1.5 by $0.47$ pp, which does not reach significance over
three seeds ($p = 0.051$), and it does not beat the ResNet-18 teacher at all
($-0.07$ pp, $p = 0.405$). Only with distillation does the margin reach $1.49$
pp and the teacher comparison reverse. Every claim in this paper about
exceeding the teacher is therefore a claim about the distilled setting.

\textbf{We cannot explain the distillation asymmetry.} Given an identical
teacher, temperature and loss weight, the proposed model gains $2.52$ pp from
distillation while ShuffleNetV2-x1.5 gains $1.50$ pp and ShuffleNetV2-x1.0
gains $1.58$ pp. The obvious explanation, that a weaker model has more room to
improve, predicts the opposite of what we observe: the proposed model has the
higher cross-entropy accuracy and still gains more. We report the asymmetry
because it is what produces the headline margin, and we decline to attribute
a mechanism to it without a test. Establishing whether this is a property of
the architecture, of the width, or of the $7\times7$ kernels would need a
study of its own.

\textbf{The accuracy is paid for in compute.} At matched parameters the
proposed model performs $2.41\times$ the operations of ShuffleNetV2-x1.5,
occupies $3.12\times$ its peak memory at batch 128, and reaches $0.66\times$
its throughput. It is cheaper than the teacher it beats on every static axis
--- $4.16\times$ fewer parameters, $2.41\times$ fewer operations, $2.60\times$
less memory at batch one --- but against the parameter-matched baseline the
exchange is accuracy for compute, and a deployment limited by arithmetic
rather than by model size should prefer the baseline.

Two further limits constrain how far the numbers travel. Both datasets are
evaluated at $32\times32$, which for Tiny-ImageNet means downsampling from the
native $64\times64$; our absolute accuracies are therefore not comparable to
published results at the native resolution, and we make no such comparison.
And while every Tiny-ImageNet configuration is run with three seeds, the
CIFAR-10 results come from a single seed each; their significance rests on
paired tests over $10{,}000$ test images within that seed, which is a weaker
guarantee than seed-level replication.

% ---------------------------------------------------------------------
\section{Conclusion}
\label{sec:conclusion}

We set out to test whether context clustering earns its place in a lightweight
classifier at low resolution, and found that it does not: at a matched
parameter budget, replacing the clustering branch with a convolutional one
improves Tiny-ImageNet accuracy by $1.03$ pp while using fewer parameters, and
the design choices that surround clustering contribute less than seed
variance. Carrying that conclusion through yields a clustering-free model
which, at 2.71M parameters with $7\times7$ depthwise kernels and logit
distillation, reaches $61.21 \pm 0.40$ on Tiny-ImageNet and $96.68$ on
CIFAR-10, exceeding a parameter-matched ShuffleNetV2-x1.5 by $1.49$ pp and its
own ResNet-18 teacher by $2.46$ pp while using $4.16\times$ fewer parameters
than that teacher. The gain is conditional on distillation and costs
$2.41\times$ the operations of the parameter-matched baseline, and we report
both alongside the measured memory and latency rather than parameter counts
alone. The ablations, configurations, seeds and per-image predictions behind
every number are released.
